Predictions that are compositions, such as softmax vectors, topic proportions or vote distributions, can be given conformal regions on the simplex, but it is unclear how such methods should be evaluated and which scores suit which conditions. We study the evaluation of split conformal regions for composition-valued responses with seven scores (Euclidean, total variation, Hellinger, KL, Linf, Aitchison, Dirichlet likelihood), in plain, locally normalised and Mondrian variants, plus CV+. We verify that the finite-sample coverage law of split conformal prediction is score-free (a Beta law, checked for three scores and two noise families), so marginal coverage cannot discriminate between scores; we show that on class labels the total-variation, Hellinger and KL scores yield exactly the LAC label sets; and we calibrate binned conditional-coverage metrics against a null simulation in which conditional coverage is exact. In a synthetic grid of $48$ configurations with a known conditional law ($100$ replicates each), the ranking of scores by region mass depends on the noise law (the Aitchison region is $2.5$ times the Euclidean one for Dirichlet noise at $K=3$ but $0.9$ times for rounded logistic-normal noise at $K=10$), local normalisation repairs conditional coverage only under heteroscedasticity and never down to the null floor, Mondrian calibration by argmax is unusable at $30$ calibration points, and Aitchison region sizes change several-fold with the zero-replacement constant. A first-order Dirichlet theory predicts pointwise Euclidean conditional coverage (RMSE $\le0.010$) but fails for the Aitchison score. Experiments on CIFAR-10 and CIFAR-10H and a covariate-shift check complete a reproducible protocol; we claim no new guarantee and limit conclusions to the tested designs.
